\ssscp{Putative PCEMP *e}{13-01-05-02}
In \citet[247--249]{Blust1993} the evidence for PCEMP *e is based on
two purported lexical developments: PMP *i-sai > PCEMP *i-sei
‘who?’ and PMP *ma-qitəm ‘black’ > PCEMP *ma-qetəm ‘black’.

Apart from these two forms (discussed further below),
we note that of all historical reconstructions listed in \citet{BlustTrussel2020}
with word-initial *e,
only one is attributed to PCEMP.
First, \citet{BlustTrussel2020} do not distinguish orthographically
between PMP *ə and PCEMP *e; both are in the database as ‹*e›.
Secondly, the one PCEMP word beginning with ‹*e› is PCEMP ‹*erit›
‘scratch, scrape’ (*erit? *ərit?), with support from only two Wallacean
languages: Ngadha \it{eri} ‘make a stroke or line,
scratch’, Kamarian \it{eri} ‘shave, scrape off’,
and four Oceanic languages beginning with /o/,
suggesting this etymon is actually *ərit, not *erit.
Thirdly, we noted in \srf{02-04-05} that many sources in the
region do not distinguish between vowel-initial
and glottal-initial words in their transcriptions.
And fourthly, PCEMP ‹*erit› ‘scratch, scrape’ should have been
linked to PMP *kərit ‘scratch,
grate’, suggesting that the few forms transcribed with an initial
vowel may actually reflect *k > \it{ʔ}, {\0}.
Not a lot of weight should be put into this etymon,
other than the observation that just as all etyma reconstructed
with initial *o are ascribed to \tsc{POc}
and not to PCEMP,
so also while there are many PMP etyma reconstructed with initial
*ə, there are no unproblematic etyma with initial *e in support of PCEMP.

With regards to putative PMP *i-sai > PCEMP *i-sei ‘who?’,
all are in agreement that this is not a replacement innovation
in the CEMP region.
Support for this development is provided from only four non-Oceanic languages.
But given the widespread tendency in many languages for the raising
of *ai > \it{e,
ei, ee} throughout the region,
a caveat is also given,
“Whether this change is a product of extensive phonologically
motivated convergence, or whether PCEMP had variants
*sai/sei is unclear” \citep[247f]{Blust1993}.

In response to arguments
and objections from \citet{Nothofer1992}
and \citet{DonohueGrimes2008}, \citet[54]{Blust2009} concedes,
“This proposed innovation [*sai > *sei] provides weaker support
for CEMP than the preceding [*oliq],
since convergence is more difficult to rule out.”

With regards to putative PMP *ma-qitəm > PCEMP *ma-qetəm ‘black’,
there is also agreement that this is not a replacement innovation
in the CEMP region.
\citet[248]{Blust1993} acknowledges that with the widespread
loss of PMP *q in the region,
“It would not be surprising if the sequence *ai in a trisyllable
had contracted independently in CMP
and Oceanic languages.” 
To put it a different way,
*-ay > \it{e}, *ai \it{> e,
ei, ee}, *aqi > \it{e, ei,
ee} occurs in many languages throughout the region with varying
combinations in different languages,
and sometimes with lexically specific variations.
So finding many languages with \it{e} in the penultimate syllable
for PMP *ma-qitəm ‘black’ is not a surprise for anyone familiar with the region.
Cognates with \it{e} are found in the Philippines and in Sulawesi.
Similar patterns are found with etyma such as *taqi ‘faeces,
excreta’, *ma-iRaq ‘red-brown’, and others.

In response to arguments
and objections from \citet{Nothofer1992}
and \citet{DonohueGrimes2008}, \citet[56]{Blust2009} says,
“but the matter remains open to question.
In conclusion, then, it probably is best to treat this comparison
[PMP *ma-qitəm > PCEMP *ma-qetəm ‘black’] with caution,
and so not include it as evidence for CEMP.”